\documentclass[11pt]{article}

\usepackage[UTF8]{ctex}
\usepackage[margin=1in]{geometry}
\usepackage{setspace}
\usepackage{hyperref}
\usepackage{booktabs}
\usepackage{longtable}
\usepackage{tabularx}
\usepackage{array}
\usepackage{enumitem}
\usepackage{amsmath,amssymb}
\usepackage{tikz}
\usepackage{xcolor}
\usepackage{fancyhdr}
\usetikzlibrary{arrows.meta,positioning,shapes.geometric,shapes.misc,fit,backgrounds}

\hypersetup{
  colorlinks=true,
  linkcolor=blue!50!black,
  urlcolor=blue!50!black,
  pdftitle={Neo 抽象账户规范},
  pdfauthor={neo-abstract-account 代码仓库规范草案}
}

\pagestyle{fancy}
\setlength{\headheight}{14pt}
\fancyhf{}
\lhead{Neo 抽象账户规范}
\rhead{2026-03-09}
\cfoot{\thepage}

\setstretch{1.08}
\setlist[itemize]{leftmargin=1.2em}
\setlist[enumerate]{leftmargin=1.4em}

\newcommand{\code}[1]{\texttt{#1}}
\newcommand{\must}{\textbf{必须}}
\newcommand{\should}{\textbf{应该}}
\newcommand{\mayword}{\textbf{可以}}
\newcommand{\aasystem}{Neo 抽象账户}
\newcommand{\AAsystem}{AA}

\title{\textbf{Neo 抽象账户规范}\\
\large Neo N3 上策略门控抽象账户的正式实现导向规范}
\author{代码仓库规范草案}
\date{2026-03-09}

\begin{document}
\maketitle
\tableofcontents
\newpage

\section{摘要}
本文档规定了在 \code{neo-abstract-account} 代码仓库中实现的\aasystem{}参考系统。该系统为 Neo N3 定义了一个策略门控的抽象账户模型，其中一个共享的主合约管理确定性的每账户验证地址、基于角色和阈值的授权、兼容 EVM 的元交易、有界的链下协作草稿、可选的中继执行以及可选的恢复验证器扩展。本规范以论文风格格式编写，面向实现者、审计员、运营者和集成者。它形式化了系统架构、执行模型、数据流、角色模型、安全边界、部署假设和恢复验证器工作流。本文档中的所有图表、表格和工作流都直接在 LaTeX 中生成，以保持可重现性和可审查性。

\section{引言}
\aasystem{}是 Neo N3 上账户抽象系统的类比，其中面向用户的账户身份与单个静态签名密钥分离。系统不是为每个账户部署专用的钱包合约，而是将账户配置存储在一个共享的主合约中，并为每个逻辑账户派生一个确定性的验证地址。这种设计降低了每账户的部署成本，同时保留了稳定的地址、单一的执行表面和一致的策略层。

本代码仓库中的实现在此基础上扩展了以下功能：
\begin{itemize}
  \item 通过管理员和经理进行原生 Neo 授权，
  \item 兼容 EVM 的 EIP-712 元交易授权，
  \item 白名单、黑名单和转账限制的策略门控，
  \item 可选的中继辅助提交，
  \item 带有作用域 URL 的链下协作交易草稿，
  \item 可选的自定义验证器和恢复验证器，
  \item Neo N3 测试网上的部署和实时验证证据。
\end{itemize}

\section{一致性语言}
本文档中的关键词\must{}、\should{}和\mayword{}表示符合代码仓库当前设计的实现的要求强度。

\begin{itemize}
  \item 符合规范的实现\must{}保留主合约执行边界。
  \item 符合规范的实现\must{}拒绝绕过 \AAsystem{} 包装器入口点的原始代理签名外部执行。
  \item 符合规范的实现\should{}保留只读、签名者和操作者角色的作用域协作模型。
  \item 符合规范的实现\mayword{}通过自定义验证器或恢复验证器扩展授权，前提是核心策略门控执行模型保持完整。
\end{itemize}

\section{术语和角色}
\begin{table}[h!]
  \centering
  \caption{核心术语}
  \begin{tabularx}{\textwidth}{>{\raggedright\arraybackslash}p{0.26\textwidth}X}
    \toprule
    术语 & 含义 \\
    \midrule
    抽象账户 & 由 \code{accountId} 标识并由共享主合约管理的逻辑账户。 \\
    确定性验证地址 & 从确定性 \code{verify(accountId)} 脚本派生的 Neo 地址。 \\
    主合约 & 共享的链上执行和授权引擎。 \\
    包装器入口点 & 主合约入口点，如 \code{executeUnifiedByAddress} 或 \code{executeUnifiedByAddress}。 \\
    草稿 & 包含操作主体、交易主体、签名和元数据的链下协作记录。 \\
    中继预检 & 在实际提交之前对中继就绪负载进行服务器支持的模拟。 \\
    恢复验证器 & 实现面向社交恢复的授权模型的专用合约。 \\
    \bottomrule
  \end{tabularx}
\end{table}

\begin{table}[h!]
  \centering
  \caption{角色和作用域模型}
  \begin{tabularx}{\textwidth}{>{\raggedright\arraybackslash}p{0.18\textwidth}>{\raggedright\arraybackslash}p{0.2\textwidth}X}
    \toprule
    层级 & 角色/作用域 & 权限 \\
    \midrule
    链上 & 管理员 & 高权限治理操作、阈值更新、验证器设置和策略更改。 \\
    链上 & 经理 & 在允许的合约和转账策略范围内的日常执行。 \\
    链下 & 公开分享链接 & 只读草稿审查。 \\
    链下 & 协作者链接 & 仅签名收集；无操作者级别的变更权限。 \\
    链下 & 操作者链接 & 中继检查、收据持久化、广播类操作和作用域链接轮换。 \\
    \bottomrule
  \end{tabularx}
\end{table}

\section{设计目标和非目标}
\subsection*{目标}
\begin{itemize}
  \item 提供确定性账户寻址，无需每账户合约部署。
  \item 在单一审计的执行表面中集中授权和策略执行。
  \item 支持原生 Neo 和 EVM 发起的签名流程。
  \item 提供操作者友好的协作和中继工具，而不将核心信任转移到链下。
  \item 通过测试网部署和验证脚本保留实时操作证据。
\end{itemize}

\subsection*{非目标}
\begin{itemize}
  \item 本规范不定义跨链账户抽象标准。
  \item 本规范不声称链下草稿存储可以替代链上授权。
  \item 本规范不定义超出代码仓库实现变体的通用标准化恢复验证器 ABI。
\end{itemize}

\section{系统架构}
\subsection{高层组件视图}
\begin{figure}[h!]
\centering
\begin{tikzpicture}[
  node distance=1.8cm and 1.8cm,
  box/.style={draw, rounded corners, align=center, minimum width=2.6cm, minimum height=1cm, fill=blue!4},
  trust/.style={draw, dashed, rounded corners, inner sep=0.25cm},
  arrow/.style={-{Latex[length=2.5mm]}, thick}
]

\node[box] (user) {用户/签名者/操作者};
\node[box, right=of user] (frontend) {前端工作区};
\node[box, below=of frontend] (drafts) {草稿存储\\(localStorage / Supabase)};
\node[box, above=of frontend] (wallets) {Neo 钱包 / EVM 钱包};
\node[box, right=of frontend] (relay) {中继和\\操作者 API};
\node[box, right=of relay] (master) {UnifiedSmartWallet\\主合约};
\node[box, above=of master] (proxy) {确定性\\验证地址};
\node[box, below=of master] (target) {目标合约};
\node[box, below=of relay] (oracle) {可选预言机/\\恢复合约};

\draw[arrow] (user) -- (frontend);
\draw[arrow] (frontend) -- (wallets);
\draw[arrow] (frontend) -- (drafts);
\draw[arrow] (frontend) -- (relay);
\draw[arrow] (frontend) -- (master);
\draw[arrow] (proxy) -- (master);
\draw[arrow] (relay) -- (master);
\draw[arrow] (master) -- (target);
\draw[arrow] (oracle) -- (master);

\node[trust, fit=(user)(frontend)(wallets)(drafts)(relay)] {};
\node[trust, fit=(proxy)(master)(target)(oracle)] {};
\node[above left=0.1cm and -0.1cm of frontend] {链下边界};
\node[above left=0.1cm and -0.1cm of master] {链上边界};
\end{tikzpicture}
\caption{系统架构和信任边界}
\label{fig:architecture}
\end{figure}

\subsection{合约模块映射}
\begin{table}[h!]
  \centering
  \caption{主合约文件职责}
  \begin{tabularx}{\textwidth}{>{\raggedright\arraybackslash}p{0.38\textwidth}X}
    \toprule
    文件 & 职责 \\
    \midrule
    \code{AbstractAccount.cs} & 顶层入口点和共享集成粘合代码 \\
    \code{AbstractAccount.AccountLifecycle.cs} & 账户创建和确定性地址绑定 \\
    \code{AbstractAccount.StorageAndContext.cs} & 存储布局、执行锁、瞬态上下文 \\
    \code{AbstractAccount.ExecutionAndPermissions.cs} & 策略门控执行和目标调用 \\
    \code{AbstractAccount.MetaTx.cs} & EIP-712 授权和签名者验证 \\
    \code{AbstractAccount.Admin.cs} & 角色和阈值治理 \\
    \code{AbstractAccount.Oracle.cs} & 预言机门控的 dome 流程 \\
    \code{AbstractAccount.Upgrade.cs} & 仅部署者更新路径 \\
    \bottomrule
  \end{tabularx}
\end{table}

\section{账户模型和状态}
\subsection{确定性寻址}
对于每个逻辑账户，系统从主合约哈希和 \code{accountId} 派生确定性验证脚本。生成的脚本哈希是面向用户的 Neo 地址。不需要每用户合约部署。符合规范的实现\must{}保留验证地址将授权路由回主合约的属性。

\subsection{链上状态类别}
\begin{table}[h!]
  \centering
  \caption{主要链上状态类别}
  \begin{tabularx}{\textwidth}{>{\raggedright\arraybackslash}p{0.24\textwidth}>{\raggedright\arraybackslash}p{0.18\textwidth}X}
    \toprule
    类别 & 示例前缀 & 目的 \\
    \midrule
    管理员集合 & \code{0x01} & 序列化的管理员地址 \\
    管理员阈值 & \code{0x02} & 所需的管理员签名数 \\
    经理集合 & \code{0x03} & 序列化的经理地址 \\
    经理阈值 & \code{0x04} & 所需的经理签名数 \\
    黑名单/白名单 & \code{0x09}, \code{0x0B} & 目标策略约束 \\
    最大转账 & \code{0x0C} & 代币转账/批准的金额上限 \\
    验证器合约 & \code{0x12} & 可选的外部验证器逻辑 \\
    Dome/预言机状态 & 多个作用域前缀 & 不活动恢复和预言机门控 \\
    \bottomrule
  \end{tabularx}
\end{table}

\subsection{链下协作状态}
协作层存储面向草稿的协调数据，而不是授权真相。草稿\must{}被视为协调工件，而不是权限来源。

\begin{table}[h!]
  \centering
  \caption{链下草稿数据}
  \begin{tabularx}{\textwidth}{>{\raggedright\arraybackslash}p{0.25\textwidth}X}
    \toprule
    字段组 & 含义 \\
    \midrule
    \code{operation\_body} & 人类/编辑器意图和合约调用形状 \\
    \code{transaction\_body} & 暂存的包装交易详情 \\
    \code{signatures} & 仅追加的收集签名和批准 \\
    \code{metadata.activity} & 有界活动时间线（最新 100 条） \\
    \code{metadata.submissionReceipts} & 有界收据历史（最新 12 条） \\
    \code{share\_slug}, \code{collaboration\_slug}, \code{operator\_slug} & 读取、签名和操作角色的作用域 URL \\
    \bottomrule
  \end{tabularx}
\end{table}

\section{授权和执行模型}
\subsection{验证阶段}
有效的确定性代理见证本身是不够的。加固后，验证阶段\must{}确认顶层交易脚本通过批准的包装器入口点返回到 \AAsystem{} 合约。

\begin{figure}[h!]
\centering
\begin{tikzpicture}[
  node distance=1.3cm and 1.0cm,
  verifybox/.style={draw, rounded corners, align=center, minimum width=2.8cm, minimum height=0.9cm, fill=green!4},
  decision/.style={draw, diamond, aspect=2, align=center, fill=yellow!10, inner sep=0.1cm},
  bad/.style={draw, rounded corners, align=center, minimum width=2.4cm, minimum height=0.8cm, fill=red!8},
  arrow/.style={-{Latex[length=2.5mm]}, thick}
]
\node[verifybox] (tx) {交易到达 Neo 节点};
\node[verifybox, below=of tx] (verify) {调用 \code{verify(accountId)}};
\node[decision, below=of verify] (selfcall) {AA 自调用？};
\node[bad, right=2.5cm of selfcall] (reject1) {拒绝代理滥用};
\node[decision, below=of selfcall] (auth) {阈值/验证器有效？};
\node[bad, right=2.5cm of auth] (reject2) {拒绝未授权调用};
\node[verifybox, below=of auth] (pass) {验证成功};
\draw[arrow] (tx) -- (verify);
\draw[arrow] (verify) -- (selfcall);
\draw[arrow] (selfcall) -- node[above] {否} (reject1);
\draw[arrow] (selfcall) -- node[right] {是} (auth);
\draw[arrow] (auth) -- node[above] {否} (reject2);
\draw[arrow] (auth) -- node[right] {是} (pass);
\end{tikzpicture}
\caption{验证阶段决策流程}
\label{fig:verification}
\end{figure}

\subsection{应用阶段}
一旦验证成功，主合约在调用目标合约之前执行应用策略。验证和应用之间的这种分离是设计的基础。

\begin{figure}[h!]
\centering
\begin{tikzpicture}[
  node distance=1.25cm and 0.8cm,
  procbox/.style={draw, rounded corners, align=center, minimum width=3.0cm, minimum height=0.85cm, fill=blue!5},
  decision/.style={draw, diamond, aspect=2.2, align=center, fill=yellow!10, inner sep=0.1cm},
  bad/.style={draw, rounded corners, align=center, minimum width=2.5cm, minimum height=0.8cm, fill=red!8},
  arrow/.style={-{Latex[length=2.5mm]}, thick}
]
\node[procbox] (entry) {调用包装器入口点};
\node[decision, below=of entry] (role) {角色/阈值满足？};
\node[bad, right=2.8cm of role] (bad1) {拒绝};
\node[decision, below=of role] (policy) {白名单/黑名单/方法策略有效？};
\node[bad, right=2.8cm of policy] (bad2) {拒绝};
\node[decision, below=of policy] (limit) {转账/批准金额在限制内？};
\node[bad, right=2.8cm of limit] (bad3) {拒绝};
\node[procbox, below=of limit] (call) {调用目标合约};
\node[procbox, below=of call] (return) {返回结果};
\draw[arrow] (entry) -- (role);
\draw[arrow] (role) -- node[above] {否} (bad1);
\draw[arrow] (role) -- node[right] {是} (policy);
\draw[arrow] (policy) -- node[above] {否} (bad2);
\draw[arrow] (policy) -- node[right] {是} (limit);
\draw[arrow] (limit) -- node[above] {否} (bad3);
\draw[arrow] (limit) -- node[right] {是} (call);
\draw[arrow] (call) -- (return);
\end{tikzpicture}
\caption{应用阶段策略执行}
\label{fig:application}
\end{figure}

\section{协议工作流}
\subsection{工作流类别}
代码仓库支持三种主要工作流类别：
\begin{enumerate}
  \item 原生 Neo 钱包执行，
  \item EVM 元交易执行，
  \item 协作草稿创建、审查和提交。
\end{enumerate}

\subsection{端到端用户工作流}
\begin{figure}[h!]
\centering
\begin{tikzpicture}[
  node distance=0.9cm,
  flowbox/.style={draw, rounded corners, align=center, minimum width=8.7cm, minimum height=0.8cm, fill=blue!4},
  arrow/.style={-{Latex[length=2.2mm]}, thick}
]
\node[flowbox] (a) {加载抽象账户或派生预期的确定性地址};
\node[flowbox, below=of a] (b) {通过预设或自定义调用编写操作};
\node[flowbox, below=of b] (c) {持久化不可变草稿主体};
\node[flowbox, below=of c] (d) {收集原生和/或 EVM 签名};
\node[flowbox, below=of d] (e) {检查就绪状态、负载模式和中继预检};
\node[flowbox, below=of e] (f) {通过客户端钱包或中继提交};
\node[flowbox, below=of f] (g) {记录收据和时间线结果};
\draw[arrow] (a) -- (b);
\draw[arrow] (b) -- (c);
\draw[arrow] (c) -- (d);
\draw[arrow] (d) -- (e);
\draw[arrow] (e) -- (f);
\draw[arrow] (f) -- (g);
\end{tikzpicture}
\caption{标准用户工作流}
\label{fig:workflow}
\end{figure}

\subsection{执行路径矩阵}
\begin{table}[h!]
  \centering
  \caption{执行路径比较}
  \begin{tabularx}{\textwidth}{>{\raggedright\arraybackslash}p{0.22\textwidth}>{\raggedright\arraybackslash}p{0.18\textwidth}>{\raggedright\arraybackslash}p{0.18\textwidth}X}
    \toprule
    路径 & 签名者来源 & 提交来源 & 注释 \\
    \midrule
    客户端原生 & Neo 钱包 & 浏览器钱包 & 最简单且默认安全的路径 \\
    中继预检 & 任何准备好的负载 & 中继 API & 仅模拟，无最终广播 \\
    中继原始提交 & 预签名原始交易 & 中继 API & 仅在明确启用原始转发时 \\
    元调用中继 & EVM 类型化数据签名 & 中继签名者 & 保留 EIP-712 UX 同时执行 \AAsystem{} 策略 \\
    \bottomrule
  \end{tabularx}
\end{table}

\subsection{协作草稿流程}
协作草稿实现\must{}保留以下作用域语义：
\begin{itemize}
  \item 公开分享链接是只读的，
  \item 协作者链接仅用于签名，
  \item 操作者链接管理中继、收据和轮换操作，
  \item 操作者级别的操作可以通过签名的服务器端变更路径进行调解。
\end{itemize}

\section{数据流和状态模型}
\subsection{系统边界视图}
\begin{figure}[h!]
\centering
\begin{tikzpicture}[
  node distance=1.7cm and 1.7cm,
  box/.style={draw, rounded corners, align=center, minimum width=2.8cm, minimum height=0.95cm, fill=orange!5},
  arrow/.style={-{Latex[length=2.5mm]}, thick}
]
\node[box] (browser) {浏览器运行时};
\node[box, right=of browser] (supabase) {Supabase 草稿存储};
\node[box, below=of browser] (local) {localStorage 后备};
\node[box, right=of local] (relay) {中继/操作者 API};
\node[box, right=of relay] (chain) {Neo N3 链};

\draw[arrow] (browser) -- (supabase);
\draw[arrow] (browser) -- (local);
\draw[arrow] (browser) -- (relay);
\draw[arrow] (browser) -- (chain);
\draw[arrow] (relay) -- (chain);
\end{tikzpicture}
\caption{跨链下和链上边界的数据移动}
\label{fig:dataflow}
\end{figure}

\subsection{权限矩阵}
\begin{table}[h!]
  \centering
  \caption{按边界划分的数据权限}
  \begin{tabularx}{\textwidth}{>{\raggedright\arraybackslash}p{0.22\textwidth}>{\raggedright\arraybackslash}p{0.28\textwidth}X}
    \toprule
    边界 & 主要数据 & 权限特征 \\
    \midrule
    链上主合约 & 角色、阈值、nonce、策略、验证器状态 & 权威的、共识支持的、最终执行源 \\
    浏览器内存 & 活动工作区和钱包会话 & 临时的、用户本地的、非权威的 \\
    localStorage & 仅本地草稿和偏好设置 & 仅同浏览器便利 \\
    Supabase & 草稿主体、仅追加签名、有界活动和收据 & 带有作用域变更的协调状态 \\
    中继 API & 模拟、提交、签名的操作者变更 & 辅助边界；不得重新定义链上权限 \\
    \bottomrule
  \end{tabularx}
\end{table}

\subsection{保留和可变性}
符合规范的协作实现\should{}保留有界元数据行为：
\begin{itemize}
  \item 保留最新 100 条活动条目，
  \item 保留最新 12 条提交收据，
  \item 不可变草稿主体独立于仅追加签名历史保留。
\end{itemize}

\section{恢复验证器}
代码仓库包含用于社交恢复式授权模式的专用恢复验证器。这些验证器与主合约分离，旨在验证专门的恢复逻辑。

\begin{table}[h!]
  \centering
  \caption{Neo N3 测试网上已验证的恢复验证器部署}
  \begin{tabularx}{\textwidth}{>{\raggedright\arraybackslash}p{0.26\textwidth}>{\raggedright\arraybackslash}p{0.31\textwidth}X}
    \toprule
    合约 & 已部署哈希 & 模型 \\
    \midrule
    ArgentRecoveryVerifier & \code{0x260b204b109506140f6e20ef99d02c142d070f72} & 带有阈值恢复流程的监护人集合 \\
    SafeRecoveryVerifier & \code{0x06a7c50c2dd81f988e2e31b7fd721501008fbfa8} & 带有时间锁参数化的模块化阈值恢复 \\
    LoopringRecoveryVerifier & \code{0x3ed17f73f19a89bc36e2dd82a19fc920aa2e54c7} & 基于监护人哈希的恢复，带有可选时间锁行为 \\
    \bottomrule
  \end{tabularx}
\end{table}

本代码仓库中的恢复验证器工作流还\must{}满足以下实现约束：
\begin{itemize}
  \item 仅使用已验证的 \code{*.Fixed.cs} 恢复源进行可重现编译，
  \item 代码仓库文档或辅助脚本中不存储明文 WIF 值，
  \item Argent、Safe 和 Loopring 存在官方包脚本验证器，
  \item 已部署的哈希和 txid 记录在代码仓库验证记录中。
\end{itemize}

\section{安全考虑}
\subsection{核心不变量}
符合规范的实现\must{}保留以下不变量：
\begin{enumerate}
  \item 确定性验证地址永远不会取代主合约作为权限源；
  \item 不支持的原始代理签名外部支出被拒绝；
  \item 验证和应用检查对于有效执行都是必需的；
  \item 链下草稿存储永远不会绕过链上授权；
  \item 仅操作者操作保持比一般协作读取或签名者变更更窄。
\end{enumerate}

\subsection{代码仓库中存在的加固措施}
\begin{itemize}
  \item 未压缩元交易公钥长度验证，
  \item 带有 \code{Retry-After} 信号的中继和操作者 API 速率限制，
  \item 清理的中继错误响应，
  \item 前端部署配置中的浏览器安全标头，
  \item 项目/构建隔离以防止跨合约树的工件污染。
\end{itemize}

\subsection{操作建议}
操作者\should{}：
\begin{itemize}
  \item 在启用中继提交之前验证运行时配置，
  \item 仅在服务器端保留密钥，
  \item 将协作链接视为作用域能力，并在暴露时轮换它们，
  \item 记录所有恢复合约的部署 txid、合约哈希和验证器输出。
\end{itemize}

\section{实现和部署说明}
\subsection{运行时配置}
代码仓库区分浏览器安全的运行时提示和仅服务器的密钥。浏览器暴露的值通知 UI 功能和端点，而签名密钥和服务角色凭据保留在服务器端。

\subsection{代码仓库验证}
代码仓库验证脚本构建主合约、重新生成合约工件、运行合约测试、运行前端测试和构建、执行前端生产审计并运行 SDK 测试套件。主合约工件生成现在通过隔离的辅助程序进行，以避免与恢复构建中间产物的交互。

\subsection{恢复验证器验证的参考工作流}
\begin{enumerate}
  \item 通过隔离的辅助程序重建恢复工件，
  \item 运行本地就绪检查，
  \item 使用 \code{neo-go contract deploy --await} 部署，
  \item 记录哈希和 txid，
  \item 使用所有三个已部署合约的环境变量运行 \code{npm run testnet:validate:recovery:all}。
\end{enumerate}

\section{测试网验证总结}
当前代码仓库状态包括 Neo N3 测试网的实时证据：
\begin{itemize}
  \item 主合约验证和加固包装器执行，
  \item 恢复验证器部署和基本链上验证，
  \item Safe 和 Loopring 恢复验证器包脚本执行，
  \item 当前分支顶端的完整代码仓库验证。
\end{itemize}

\section{结论}
\aasystem{}将确定性寻址、共享合约执行、作用域链下协作、中继辅助提交和专门的恢复验证器组合成一个连贯的策略门控账户抽象模型，用于 Neo N3。该设计优先考虑明确的信任边界、低每账户部署开销和可重现的验证。代码仓库的当前实现和验证证据支持实现者和审计员将其用作操作上有根据的参考系统，而不是纯粹的理论抽象设计。

\appendix
\section{附录 A：代码仓库文档映射}
以下文件补充了本正式规范：
\begin{itemize}
  \item \code{README.md}
  \item \code{docs/INDEX.md}
  \item \code{docs/HOW\_IT\_WORKS.md}
  \item \code{docs/WORKFLOWS.md}
  \item \code{docs/DATA\_FLOW.md}
  \item \code{contracts/recovery/TESTNET\_VALIDATION\_2026-03-09.md}
\end{itemize}

\section{附录 B：推荐读者配置}
\begin{itemize}
  \item \textbf{操作者}：重点关注第 6、7、8 和 9 节。
  \item \textbf{审计员}：重点关注第 5、6、8 节和恢复验证器附录。
  \item \textbf{集成者}：重点关注第 4、6、7 和 9 节。
  \item \textbf{协议工程师}：按顺序阅读文档。
\end{itemize}

\end{document}
